\begin{table}[H]
\centering
\small
\setlength{\tabcolsep}{4pt}
\begin{tabularx}{\textwidth}{@{}>{\raggedright\arraybackslash}X*{5}{>{\centering\arraybackslash}X}@{}}
\toprule
Model & \makecell{Intra val\\PR} & \makecell{Intra test\\PR} & \makecell{Intra test\\ROC} & \makecell{BIOGRID\\PR} & \makecell{BIOGRID\\ROC} \\
\midrule
Seq LogReg & 0.60 & 0.62 & 0.62 & 0.76 & 0.78 \\
Seq MLP & 0.67 & 0.68 & 0.70 & 0.78 & 0.79 \\
Seq+Text MLP & 0.68 & 0.70 & 0.71 & 0.78 & 0.80 \\
Seq+Text Attn & 0.69 & 0.70 & 0.71 & 0.79 & 0.80 \\
Attn Ensemble & 0.70 & 0.72 & 0.72 & 0.80 & 0.81 \\
Log-loss Ens. & 0.64 & 0.65 & 0.67 & 0.80 & 0.82 \\
\bottomrule
\end{tabularx}
\caption[Predictive metrics for the frozen surrogate family]{Predictive benchmarking table for the frozen surrogate family used in the protein application. The original attention ensemble remains the final thesis surrogate; the log-loss-optimized ensemble is shown as a calibration-oriented contrast because it improves loss-based metrics but does not pass the balanced replacement rule. Intra val denotes the in-domain validation split used for model selection and later reused as the pilot/calibration split in the Chapter~5 planning calculations.}
\label{tab:chap5_predictive_metrics}
\end{table}
